\documentclass[border=4pt]{standalone}
\usepackage[siunitx]{circuitikz}
\usetikzlibrary{calc}

\begin{document}
\begin{circuitikz}[font=\small,line width=0.9pt]
  \ctikzset{tripoles/thickness=1,logic ports/thickness=1,
    flipflops/thickness=1}
  \foreach \i in {0,...,3} {
    \pgfmathsetmacro{\mx}{1.15+3.35*\i}
    \pgfmathsetmacro{\fx}{2.83+3.35*\i}

    % Conventional two-input multiplexer: wide input side, tapered output side.
    \draw (\mx-0.48,4.42) -- (\mx+0.40,4.15) --
      (\mx+0.40,3.55) -- (\mx-0.48,3.30) -- cycle;
    \coordinate (MP\i) at (\mx-0.48,4.20);
    \coordinate (MS\i) at (\mx-0.48,3.48);
    % Exact intersection of x=mx with the lower MUX edge from
    % (mx-0.48,3.30) to (mx+0.40,3.55).
    \coordinate (MC\i) at (\mx,3.43636);
    \node at (\mx-0.28,4.20) {1};
    \node at (\mx-0.28,3.48) {0};
    \node[font=\scriptsize] at (\mx+0.05,3.85) {MUX};

    \node[flipflop D,scale=0.78] (F\i) at (\fx,3.40) {};
    \coordinate (MO\i) at (\mx+0.40,0 |- F\i.pin 1);
    \node[above] at (F\i.north) {stage $\i$};
    \draw (MO\i) -- (F\i.pin 1);
    \pgfmathsetmacro{\px}{\mx-0.82}
    \draw (\px,4.85) node[above] {$P_{\i}$} |- (MP\i);
    \draw (MC\i) -- (\mx,2.92) -- (\mx-0.42,2.92)
      node[left] {$L$};
  }

  \draw (0.05,3.48) node[left] {$SER_{in}$} -- (MS0);
  \foreach \i/\j in {0/1,1/2,2/3} {
    \coordinate (QS\i) at ($(MS\j)+(-0.45,0)$);
    \draw (F\i.pin 6) -- (F\i.pin 6 -| QS\i)
      -- (QS\i) -- (MS\j);
  }
  \draw (F3.pin 6) -- ++(0.55,0)
    node[right] {$SER_{out}$};

  \draw (0.05,1.05) node[left] {$CLK$} -- (13.60,1.05);
  \foreach \i in {0,...,3} {
    \coordinate (FC\i) at ($(F\i.pin 3)+(-0.10,0)$);
    \coordinate (FB\i) at (FC\i |- 0,1.05);
    \draw (FB\i) -- (FC\i) -- (F\i.pin 3);
    \fill (FB\i) circle (1.4pt);
  }
\end{circuitikz}
\end{document}
